\documentclass[11pt]{article}
\usepackage[utf8]{inputenc}
\usepackage[T1]{fontenc}
\usepackage[spanish]{babel}
\usepackage[margin=2.5cm]{geometry}
\usepackage{amsmath,amssymb}
\usepackage{enumitem}
\setlength{\parindent}{0pt}
\setlength{\parskip}{6pt}

\begin{document}

\begin{center}
{\Large\bfseries Microeconomía II --- Control 1 (práctica)}\\[4pt]
{\large Discriminación de 1.\textsuperscript{er} grado y tarifa en dos partes}\\[4pt]
20 puntos \quad$\cdot$\quad 20 minutos \quad$\cdot$\quad libro cerrado
\end{center}

\hrule
\medskip

Un gimnasio es el único del pueblo. La demanda de un socio típico por clases al mes es
\[
q = 60 - P ,
\]
y el costo marginal de cada clase es constante e igual a $CMg = 10$. No hay costo fijo.

\begin{enumerate}[label=\alph*), leftmargin=*, itemsep=10pt]

\item \textbf{Precio único (5 pts).} Si el gimnasio cobra el mismo precio por cada clase, calcule $Q$, $P$ y $\pi$. Calcule además el excedente del consumidor ($EC$) y la pérdida irrecuperable de eficiencia ($PIE$) respecto de competencia perfecta.

\item \textbf{Discriminación de 1.\textsuperscript{er} grado (5 pts).} Suponga que el gimnasio conoce la disposición a pagar por cada clase y cobra exactamente eso. Calcule la cantidad, el precio de la última clase vendida, el ingreso total, $\pi$, $EC$ y $PIE$. Explique en una línea, usando el ingreso marginal, por qué la cantidad coincide con la de competencia perfecta.

\item \textbf{Tarifa en dos partes (6 pts).} Ahora el gimnasio cobra una cuota mensual $A$ más un precio $p$ por clase.
\begin{enumerate}[label=\roman*., itemsep=4pt]
  \item Determine $p^{*}$ y $A^{*}$, y calcule $\pi$.
  \item Compare $\pi$ con los resultados de a) y b).
  \item Explique por qué, si todos los socios son idénticos, al gimnasio no le conviene fijar $p$ por encima del $CMg$.
\end{enumerate}

\item \textbf{Comente (4 pts).} «La discriminación de 1.\textsuperscript{er} grado es eficiente; por lo tanto, es socialmente deseable.» Indique si es verdadero, falso o incierto y justifique en 3--4 líneas.

\end{enumerate}

\medskip
\hrule
\smallskip
{\small\textit{Indicaciones:} muestre el desarrollo completo (se corrige el procedimiento), trabaje con fracciones, grafique a) y b) en los mismos ejes y cierre cada parte con una conclusión.}

\end{document}
